\section{테스트 벤치 사양: 재현하는 법}
\label{sec:dishspec}

아래에는 이런 테스트 벤치를 다시 조립하는 데 필요한 모든 것을 모았다. 장비, 루프, 입력 부호화, 출력 판독, 되먹임, 실험 프로토콜이다. 매개변수마다 출처를 표시했다. “논문”은 값을 논문~\cite{Kagan2022}의 본문에서 가져왔다는 뜻이다. “선택”은 논문에 값이 없어서 재현하려면 직접 정해야 한다는 뜻이다. 이때는 기본값을 제안하고 그것을 확인하는 방법을 적는다.

\subsection*{장비와 배양}

\begin{center}
\footnotesize
\setlength{\tabcolsep}{4pt}
\renewcommand{\arraystretch}{1.15}
\begin{tabular}{>{\raggedright}p{4.0cm}>{\raggedright}p{6.9cm}>{\raggedright\arraybackslash}p{1.6cm}}
\toprule
매개변수 & 값 & 출처 \\
\midrule
칩 & MaxOne 배열: $8$~mm$^2$ 면적에 백금 전극 $26\,000$개 & 논문 \\
기록 & 동시에 최대 $1024$채널, 초당 $20\,000$ 표본 & 논문 \\
자극 & 기술적으로는 최대 $32$개 전극, 실험에서는 독립 전극 $8$개 & 논문 \\
세포 & 생쥐 태아 피질. 줄기세포에서 얻은 인간 뉴런(두 가지 배양법). 대조용 HEK293T 세포(뉴런 아님) & 논문 \\
파종 & 칩당 약 $10^6$개 세포 & 논문 \\
실험 시점의 배양 나이 & 생쥐: $\sim10$일부터. 인간: 배양법에 따라 $14$--$24$일 또는 $\sim80$일 & 논문 \\
\bottomrule
\end{tabular}
\end{center}

\subsection*{루프와 시간}

루프는 $10\ms$ 창($200$ 표본) 단위로 돈다. 창마다 운동 영역 전극의 스파이크를 세고, 게임을 갱신하고, 다음 자극을 내보낸다. 배양 속 스파이크에서 응답 자극까지의 지연은 약 $5\ms$이다. 게임 중 각 전극의 신호는 차단 주파수 $100$~Hz인 2차 베셀 고역 통과 필터를 지난다. 신호의 절댓값은 $1$~Hz의 1차 베셀 저역 통과 필터로 평활하며, 스파이크 문턱값은 이 평활된 절댓값에 비례한다. 배경의 표준편차 여섯 배라는 문턱값은 논문이 게임이 아니라 개별 활동 측정에 대해 밝힌 것이다. 자극이 배양의 응답으로 세어지지 않도록, $75$~mV를 넘는 큰 스파이크가 동시에 $15$개보다 많이 오면 모든 신호를 차단한다. 모두 논문에 따른 것이다.

\subsection*{공의 부호화}

\begin{center}
\footnotesize
\setlength{\tabcolsep}{4pt}
\renewcommand{\arraystretch}{1.15}
\begin{tabular}{>{\raggedright}p{3.4cm}>{\raggedright}p{7.5cm}>{\raggedright\arraybackslash}p{1.6cm}}
\toprule
매개변수 & 값 & 출처 \\
\midrule
감각 영역 & 전극 $8$개. 이웃한 전극이 이웃한 공 위치를 뜻하도록 배치 & 논문 \\
위치 부호 & 전극 번호 $k=1,\dots,8$이 라켓 중심에 대한 공의 위치를 정한다 & 논문 \\
어느 좌표인가 & 공의 수직 위치. 재현에서는 라켓 기준으로, $y-p\in[-\tfrac12,\tfrac12]$일 때 $k=\lceil 8(y-p+\tfrac12)\rceil$ & 논문; 공식은 선택 \\
발화율 부호 & 초당 $4$--$40$회의 자극 빈도가 라켓까지의 거리를 전한다 & 논문 \\
눈금 방향 & 공이 반대쪽 벽에 있으면 초당 $4$회, 라켓 쪽 벽에 가까워질수록 선형으로 초당 $40$회까지: $f=4+36\,(1-x/L)$, $x$는 라켓 쪽 벽까지의 거리, $L$은 코트 너비 & 논문 \\
펄스 모양 & 짧은 직사각형 2상 펄스: 먼저 양의 위상, 다음 음의 위상 & 논문 \\
진폭 & $75$~mV. 억제된 세포를 억지로 발화시키지 않도록 고른 값 & 논문 \\
위상 길이 & 밝히지 않음. 자극 시스템의 표준 설정을 쓰고 실험 중에는 바꾸지 않는다 & 선택 \\
\bottomrule
\end{tabular}
\end{center}

\subsection*{라켓 판독}

논문에는 운동 영역이 두 개 있다. 첫째 영역의 스파이크는 라켓을 위로, 둘째 영역의 스파이크는 아래로 움직인다. 세포 밀도와 활동의 차이를 반영하도록 두 영역의 기여에 가중치를 준다~\cite{Kagan2022}. 정확한 공식은 나와 있지 않다. 재현에서는 $10\ms$ 창마다 규칙~\eqref{eq:dishreadout} $\Delta p=c\,(n_\uparrow-n_\downarrow)$를 쓰고, 휴지 상태의 평균 활동으로 두 영역을 맞추는 가중치를 준다(선택). 즉 $n_\uparrow$와 $n_\downarrow$를 게임 전에 잰 각 영역의 창당 평균 스파이크 수로 나눈다. 계수 $c$는 평균 스파이크 수 하나만큼의 차이가 창마다 라켓을 코트 높이의 $1\%$만큼 옮기도록 고른다(선택). 그러면 한 영역이 계속 우세할 때 라켓은 $1$~s 동안 코트 전체를 가로지른다.

칩 위의 영역 배치는 논문 본문에 좌표로 주어지지 않았고 모식도만 있다. 재현에는 겹치지 않는 전극 그룹 세 개면 충분하다. 가운데에 감각 전극 여덟 개, 양옆에 기록 전극 그룹 하나씩, 하나는 “위”, 다른 하나는 “아래”이다(선택).

\subsection*{게임}

코트 크기, 라켓 길이, 공의 속도는 논문에 없다. 공이 일정한 속도로 날아가 벽에 튕긴다는 말뿐이다. 재현용(선택): 코트는 $1\times1$, 라켓은 오른쪽 벽에 길이 $0.2$(모델~\texttt{models/dishbrain\_model.py}에서처럼 $|y-p|<0.1$이면 히트), 공은 $1$~s 만에 코트를 가로지른다. 한 번도 받아 치지 못하고 놓친 랠리는 “에이스”, 연속 세 번보다 많이 받아 친 랠리는 “긴 랠리”로 센다(논문).

\subsection*{되먹임}

\begin{center}
\footnotesize
\setlength{\tabcolsep}{4pt}
\renewcommand{\arraystretch}{1.15}
\begin{tabular}{>{\raggedright}p{2.6cm}>{\raggedright}p{8.3cm}>{\raggedright\arraybackslash}p{1.6cm}}
\toprule
사건 & 주는 것 & 출처 \\
\midrule
히트 & 예측 가능한 자극: 감각 전극 $8$개 모두에 동시에, $75$~mV, 초당 $100$회, $100\ms$. 그동안 평소의 공 자극은 중단된다 & 논문 \\
미스 & 예측 불가능한 자극: $150$~mV, 초당 $5$회, $4$~s, 무작위 시점에 $8$개 가운데 무작위로 고른 전극에. 그 뒤 $4$~s 동안 자극 없음, 공은 무작위 방향으로 출발 & 논문 \\
무작위성의 법칙 & 밝히지 않음. 재현에서는 펄스 시점을 평균 빈도 초당 $5$회의 푸아송 과정으로 & 선택 \\
\bottomrule
\end{tabular}
\end{center}

\subsection*{프로토콜과 대조 실험}

세션 하나는 $20$~분이며, 처음 $5$~분과 마지막 $15$~분을 비교한다(논문). 결과 척도는 세 가지이다. 평균 랠리 길이, 에이스 비율, 긴 랠리 비율. 실험 조건(논문):
\begin{itemize}
\item \emph{자극}: 위에서 설명한 완전한 루프.
\item \emph{침묵}: 히트나 미스의 자극 대신 같은 시간 동안 자극을 전혀 주지 않고, 그 뒤 공이 무작위 방향으로 출발한다.
\item \emph{되먹임 없음}: 미스 뒤에 공은 그냥 튕겨 계속 날아가고, 공 위치 자극은 이어진다.
\item \emph{휴지}: 게임은 진행되고 라켓은 활동에 따라 움직이지만 자극은 전혀 없다.
\item \emph{대조}: 배지만 있는 칩. HEK293T 세포. 세포 대신 모의실험을 쓴 “컴퓨터 속 배양”.
\end{itemize}
본 실험의 표본 크기: 생쥐 배양 $9$개(세션 $101$번), 인간 배양 $11$개($138$), 배지만 있는 칩 $6$개($80$), 휴지 배양 $20$개($42$), 모의실험 $3$개($38$), 모두 $399$세션. 되먹임 실험에서는 $3$일 동안 하루 $3$세션씩 $486$세션, 조건은 “자극”, “침묵”, “되먹임 없음”($n=27$, $17$, $15$). 처음 $5$분에서 마지막 $15$분으로 평균 랠리 길이가 늘어난 것은 살아 있는 배양에서만이다. 되먹임 실험에서 시간에 따른 유의미한 향상은 “자극” 조건에서만 있었다. 세션 끝에 “침묵” 조건은 “되먹임 없음”보다 여전히 나았지만 효과가 작았고, 3일 동안 세션 간의 안정적인 학습은 없었다~\cite{Kagan2022}.

\subsection*{테스트 벤치의 주기, 단계별로}

테스트 벤치 컴퓨터는 $10\ms$마다 다음을 한다.
\begin{enumerate}
\item 지난 창 동안 두 운동 영역의 스파이크 수 $n_\uparrow$, $n_\downarrow$를 읽고, 각 영역의 휴지 상태 평균 스파이크 수로 정규화한다.
\item 라켓을 $\Delta p=c\,(n_\uparrow-n_\downarrow)$만큼 옮기고 코트 가장자리 안으로 제한한다.
\item 공을 한 걸음 옮기고, 위쪽, 아래쪽, 왼쪽 벽에서 튕긴다.
\item 공이 오른쪽 벽에 닿았으면: $|y-p|<0.1$일 때 히트, 랠리 수가 하나 늘고 히트 자극($100\ms$)을 준다. 아니면 미스, 랠리를 기록하고 미스 자극($4$~s)을 준 뒤 $4$~s 쉬고 공이 다시 출발한다.
\item 그렇지 않으면 공의 수직 위치로 전극 $k$를, 라켓 쪽 벽까지의 거리로 빈도 $f$를 골라 다음 창까지 전극 $k$에 빈도 $f$로 $75$~mV 2상 펄스를 준다.
\end{enumerate}
이것만으로 발표된 것과 호환되는 테스트 벤치를 조립하고 이 장의 핵심 질문을 확인할 수 있다. 배양을 가르치는 것이 예측 가능성인가, 아니면 히트와 미스에 대한 응답의 단순한 차이인가? 그러려면 논문의 세 조건에 논문에 없는 네 번째 조건을 더하는 것이 좋다. 미스 뒤에 예측 불가능한 자극과 같은 길이, 같은 에너지의 \emph{예측 가능한} 자극을 주는 것이다. 그래도 배양이 학습한다면 중요한 것은 예측 가능성이 아니라 차이이다.
